\section{Threshold sets and local coordinates}
\label{sets:section}

We construct cuts by assigning a measurable set of labels to every point
of the column model. Along a column these sets will grow with height,
with symmetric-difference measure equal to the vertical length. At a
switch the sets will agree exactly. Every resulting cut distance is
therefore bounded by the length of every column path. The local updates
must produce lower separation while preserving these two properties.

There are two systems, called $A$ and $B$. System $A$ also uses the
connected components of a label's active columns to separate points.
System $B$ supplies a second local separation test. Evolving distributions
of monotone cuts through smaller scales is a central mechanism in
Sidiropoulos's nonpositively curved planar construction~\cite[Sections~4--5]{Sidiropoulos13}.
Here the updates must also preserve the exact sets where local formulas meet.

\subsection{Charts on a fixed label space}

For each system fix a label space
\[
 (\Omega,\mu)=\bigl(\R\times\{-1,1\}^{\N},
                   \mathcal L\otimes\nu\bigr),
\]
where $\mathcal L$ is Lebesgue measure and $\nu$ is the product of fair sign
measures. Every set in that system is a subset of this same $\Omega$.
The random choices of updates are independent of the label coordinates;
we condition on those choices when describing the sets.

A \emph{chart} consists of a region $\mathcal R\subset X$, a height band
$I$ containing its heights, a field $P$ on the corresponding component of
$X_I$ (or on all of $X_I$), and measurable functions
$Z:\Omega\to\R$ and $\sigma:\Omega\to\{-1,1\}$ whose joint pushforward
measure is Lebesgue measure times a fair sign. For $v\in\mathcal R$ set
\begin{equation}\label{sets:formula}
 \mathcal A_v=\{\omega\in\Omega:
 Z(\omega)<h(v)+\sigma(\omega)\kappa P(v)\}.
\end{equation}
Write $\mathcal B_v$ for system $B$. The positive constant $\kappa$ is
fixed separately in each system; its values $\kappa_A,\kappa_B$ will be
chosen in Section~\ref{sets:parameters}.

At each stage there are finitely many chart identifiers. Each identifies
common field and label data and their current region. These regions
partition $X$, with a fixed rule for ties. A region may be disconnected;
its connected pieces do not receive different identifiers. When some
points receive a new chart, the old identifier is retained on the
remaining region. Initially each system has one chart on $X$, with
$P=0$, $Z$ the original Lebesgue coordinate, and $\sigma$ the first sign
coordinate.

Write $\Lip_{\mathrm{vert}}(P)$ for the supremum of the Lipschitz constants
of $P$ on its vertical column intervals. The following identity explains
both the upper invariant and the purpose of changing the field.

\begin{lemma}[Vertical control and rebasing]\label{sets:vertical}
For two points whose sets are expressed in common chart variables,
\begin{equation}\label{sets:distance}
 \mu(\mathcal A_v\mathbin\triangle\mathcal A_w)
   =\max\{|h(v)-h(w)|,\kappa|P(v)-P(w)|\}.
\end{equation}
If $\kappa\Lip_{\mathrm{vert}}(P)<1$ and $P$ respects switches, the sets are
nested increasingly on each column interval in the chart. Their
symmetric-difference measure there equals the vertical length.

A branch $P(v)=\psi(h(v))$, where $\psi:\R\to\R$ is Lipschitz with
$\kappa\Lip(\psi)<1$, can be represented by a fresh constant chart without
changing any sets.
\end{lemma}
\begin{proof}
For each sign, the symmetric difference of the two half-lines has measure
$|\Delta h+\sigma\kappa\Delta P|$. Averaging the signs gives
\eqref{sets:distance}. The functions $H\mapsto H\pm\kappa P(i,H)$
are increasing, and their average increment is the height increment.

For rebasing put $\phi_\sigma(H)=H+\sigma\kappa\psi(H)$.
These are increasing bijections of $\R$. Define the new coordinate on the
same label space by $Z'(\omega)=\phi_{\sigma(\omega)}^{-1}(Z(\omega))$.
Its pushforward measure is Lebesgue measure, since for $a<b$,
\[
 \mu\{a<Z'<b\}
 =\frac{\phi_+(b)-\phi_+(a)+\phi_-(b)-\phi_-(a)}2=b-a.
\]
Choose a sign coordinate unused anywhere earlier in this execution. It
is independent of $Z'$. Thus $(Z',\sigma')$ has the required joint
pushforward, and, for every label,
\[
 \{Z<\phi_\sigma(h)\}=\{Z'<h\}.
\]
This describes exactly the same sets by a chart with field zero.
\end{proof}

In an \emph{ordinary chart}, the field has the form
$P=c+\eta s$, where $\eta\in\{-1,1\}$ and either
$s(v)=d_I(o,v)$ for a fixed source $o$, or $s=0$.
For a constant field take $c=P$ and $\eta=1$.
Updating in \emph{ordinary orientation} means changing $s$ and converting
back using the same $c$ and $\eta$. System $A$ always uses ordinary charts.
System $B$ begins with an ordinary chart and later uses the additional
field type constructed in Section~\ref{ord:section}. Every recorded
$B$-field will have Lipschitz constant less than $5$ for finite $d_I$.

The interpolation in Section~\ref{int:section} will produce formulas for
ordinary charts. These formulas carry names: when the same formula is inherited
across an interpolation boundary, its name is preserved. When an interpolation formula becomes a chart, its new identifier
records the old chart, the particular interpolation, and the selected
formula's name, and retains the old label functions. Different old charts are not merged merely
because their formulas agree. Within one old chart, recording an inherited
formula therefore does not split a region merely because an internal
interpolation boundary was crossed.

\subsection{Exact gluing and component refinement}

Every local update is first expressed in the old chart variables. A
clipping replaces the old comparison field by its maximum with a proposal,
which is the minimum of finitely many fields. Several proposals are
combined by a finite maximum. On the boundary of each proposal's
domain the proposal lies below the old field, with a strict margin in
the spatial constructions below. Thus the clipping equals the old field
at the endpoints of every updated column interval and at any switch
leading outside the domain. Inside the domain all branches respect
switches and have vertical Lipschitz constants whose products with
$\kappa$ are less than one.

The new sets consequently retain the global upper invariant. For points
in one updated interval this follows from Lemma~\ref{sets:vertical}. For
points in different updated intervals, compare the first point with its
interval endpoint, then use the old nested sets between the two interval
endpoints, and finally compare with the second point. These three
comparisons preserve nesting, and their symmetric-difference measures
sum to the full vertical length. Unchanged points are handled by the
same comparison. This argument does not require a finite number of
updated intervals. Exact agreement at the boundary also preserves the
sets at every switching identification. A Lipschitz perturbation that
vanishes at the domain boundary admits the same interval-endpoint
argument, provided the updated field respects switches and still has
$\kappa\Lip_{\mathrm{vert}}(P)<1$. The later increments in system $B$
will satisfy these conditions.

Only after making the update do we record its winning formulas as new
charts. A spatial branch inherits the old label functions $Z,\sigma$;
a height-only branch is rebased as in Lemma~\ref{sets:vertical}. An
unchanged point keeps its old identifier. Height branches can have larger
slopes than the recorded ordinary or $B$-fields; after rebasing their
recorded fields are zero. Recording or rebasing these charts does not further change the newly
assigned subsets of $\Omega$. The identity \eqref{sets:distance} applies when the sets under
comparison are written in common label variables; the global nesting
and vertical-length invariant applies across all charts.

For a fixed label $\omega$, call a point \emph{active} when its assigned
set contains $\omega$. The active points of any column form an upper
interval. Besides membership, we can use the connected components joined
by active switches.

\begin{lemma}[Component refinement]\label{sets:component-refinement}
Suppose a system of sets is nested increasingly on each column, respects
switches, and has symmetric-difference measure at most vertical length.
It defines an $L_1$ pseudometric dominated by $d$. It also admits the
following refinement with the same upper bound.

For a label $\omega$, join two columns when one of their switches occurs
at a height active for $\omega$. To an active point assign the unit vector
indexed by its component in this finite graph, and to an inactive point
assign zero. Integrating the $\ell_1$ distances of these vectors gives the
refined pseudometric. It dominates the original symmetric-difference
pseudometric.

If a non-neighbor link has inactive tip and two active points have
representatives strictly on opposite sides of it, their refined distance
for that label is $2$.
\end{lemma}
\begin{proof}
For a fixed label all active points of a column have the same component
coordinate. On a vertical segment the vector changes only when activity
changes, and its $\ell_1$ cost is the original indicator cost. At an active
switch the two columns belong to the same component; at an inactive switch
both vectors are zero. Integration and any finite path give the upper
bound. A membership difference always costs one, proving domination.

Every switch incident to an endpoint column of a non-neighbor link is at
or below its tip. If the tip is inactive, nesting makes all such switches
inactive. A path of column indices between the two strict sides must pass
through an endpoint of the link, by noncrossing. No active path can do so.
The active points therefore have different unit coordinates.

There are only finitely many possible components, which can be indexed by
subsets of the columns. An active link can be tested at its maximum
switching height, so all coordinates are measurable. Subtracting the vector
assigned to one fixed reference point makes each assignment integrable:
the integrated norm of the difference is bounded by the distance to that
point. The same subtraction works for the unrefined indicator system.
\end{proof}

Only system $A$ will be refined, once, after the last scale. We therefore
track changes of membership during the construction, without tracking
changes of component coordinates. The final component separation will
be certified by activity at two vertices and inactivity at one fixed tip.

\subsection{The estimates required of a local update}

An update at scale $r>0$ will try to separate original vertex pairs
whose distances lie in $[r,2r)$. At scale $r$, a pair $(x,y)$ is
\emph{eligible} if the two closed balls $B_d(x,5r)$ and $B_d(y,5r)$
lie in one common old chart region in system $A$ and in one common
old chart region in system $B$. The common region may differ between
the two systems. This condition supplies common field and label data
for all short paths used in the pair's local tests.

The construction below will give constants $M,C_0<\infty$ and
$b,p,\rho_0>0$, independent of the graph and the update history,
with three properties.
\begin{enumerate}
\item At every point of the full column space $X$, an update at scale
$r$ changes each unrefined label set by symmetric-difference measure
at most $Mr$.
\item For a vertex pair with $d(x,y)\in[r,2r)$, conditional on any past
in which $(x,y)$ is eligible, the update
has probability at least $p$ of producing either unrefined symmetric-difference separation
at least $br$ in one system, or a non-neighbor link separating
representatives of $x,y$ strictly and a system-$A$ label set of measure
at least $br$ on which $x,y$ are active and the link's tip is inactive.
\item Suppose a fixed ball $B_d(v,\rho r)$ lies in one old chart region in each
system, where $0<\rho\le\rho_0$. Conditional on the past, the probability
that this ball meets different new chart regions in either system is
at most $C_0(\rho+\eta_0)$.
\end{enumerate}
Here $\eta_0$ is a dimensionless upper bound for the relative mesh of
the finite grids used to sample update parameters. It can be chosen
arbitrarily small after the finite set of scales is fixed; the
sampling convention is given in Section~\ref{sets:setup}.

Propositions~\ref{pad:success} and~\ref{pad:persistence} establish
these estimates after the local construction is complete. The third
estimate keeps a pair eligible until its own scale with positive
probability. The first then prevents sufficiently smaller scales from
erasing the second estimate's gain. For a component certificate, this
requires movement control at the tip as well as at the two vertices;
that is why the first estimate concerns all of $X$. Section~\ref{pad:assembly}
combines the estimates by using schedules whose successive scales
differ by a sufficiently large fixed factor.
